\section*{Synthèse pour la direction}
\addcontentsline{toc}{section}{Synthèse pour la direction}

MEDICA-SERVICES héberge et exploite les applications de gestion de douze établissements de soins privés \piece{(P01)}. Ses environnements portent des identifiants de patients et des données de santé \piece{(P05)}. Ces enjeux fixent son \textbf{niveau adéquat de maturité à \Cible}, c'est-à-dire le niveau de rigueur qu'elle devrait viser.

Les pièces démontrent un \textbf{niveau 1} pour la stratégie, les risques, la supervision et la réalisation des mesures, et un \textbf{niveau 2} pour les règles, la conception et l'exploitation. Aucun processus n'atteint le niveau 3, d'où un écart de deux à trois niveaux. L'affirmation de la DSI selon laquelle tous les comptes privilégiés ont un MFA est contredite, puisque seuls six sur dix sont enrôlés \piece{(T04)}. Celle de la DG sur les tests trimestriels de restauration n'est pas démontrée, deux services sur neuf ayant un procès-verbal \piece{(T09)}.

Trois risques dominent. Le premier est la prise de contrôle d'un compte à privilèges, illustrée par la connexion d'admin-old le 9 octobre suivie d'un ajout à un groupe d'administrateurs, puis par la fermeture automatique de l'alerte AL-1048 sans examen tracé \piece{(T07, T08)}. Ces signaux appellent une investigation sans prouver une intrusion. Le deuxième est l'impossibilité de reprendre le service en 8 heures après un rançongiciel, car les postes peuvent joindre les sauvegardes \piece{(FW-31)}. Le troisième concerne APP01, à qui manque un correctif critique \piece{(T05, T06)}.

Nous demandons trois décisions à la direction. Il s'agit d'abord de lancer dès maintenant l'investigation d'AL-1048 et les mesures conservatoires. Il s'agit ensuite d'engager \TotalRetenu\ sur \Plafond\ pour six actions. Il s'agit enfin d'accepter par écrit, pour douze mois, les risques résiduels des mesures différées.
